% Standalone article fragment; parent article supplies theorem environments.
% Suggested references: hatcher-covering, aht-orbits, lackenby2026.
\section{Symbolic gluing of regular covering blocks}
\label{sec:regular-cover-gluing}

\subsection{The new interface and its exact scope}

The existing surface-cover kernel records component types and transports
ordered fibre points within one supplied cover.  We now give an exact
assembly interface: it compares multiple covering blocks after their
boundary preimages have been glued, preserves named boundary lifts, and
returns a complete family of compatible transports.  In the unmarked case
it also constructs a canonical cache key.  The computational parameter is
the binary length of the covering degree; no sheet or boundary-lift list
is expanded.

The ingredients from covering-space theory are classical
\cite[Section 1.3]{hatcher-covering}.  Spanning-tree elimination,
simultaneous conjugacy, generalized Chinese remaindering, and Burnside
counting are also standard algebraic tools.  Our contribution is their
precise combination into a checked geometric representation and a
proof-carrying implementation for this repository.  We make no claim of
worldwide novelty for the underlying classification of covering spaces.

Fix one of the groups
\[
 C_m=\mathbb Z/m\mathbb Z,\qquad
 D_m=\langle T,R\mid T^m=R^2=1,\ RTR=T^{-1}\rangle.
\]
Throughout, $D_m$ has \emph{order $2m$}.  Its elements are represented by
$(s,t)=T^tR^s$, where $s\in\{0,1\}$ and $t\in\mathbb Z/m\mathbb Z$, with
\[
 (s,t)(u,v)=(s+u,\ t+(-1)^s v),\qquad
 (s,t)^{-1}=(s,-(-1)^s t).
\]
For $C_m$ the parity is always zero.  The cases $m=1,2$ are included:
$D_1\cong C_2$ and $D_2\cong C_2\times C_2$.

Let $\Gamma$ be a finite labelled multigraph; loops and parallel edges are
allowed.  At vertex $v$, take a compact connected orientable bordered
surface $F_v$ and a surjective homomorphism
\[
 \rho_v:\pi_1(F_v)\longrightarrow G,\qquad G=C_m\text{ or }D_m.
\]
The corresponding block cover $\widetilde F_v\to F_v$ has fibre $G$ and
monodromy acting on the \emph{left}.  Surjectivity makes the block cover
connected and regular.  Each boundary circle has a based peripheral word,
including a chosen path from the block basepoint, hence a specified
monodromy element $p_{v,j}\in G$.  These paths are part of the framing.

An oriented edge $e:u\to v$ uses one distinct boundary port at each end,
with
\[
 p_{v,j}=p_{u,i}^{-1}.
\]
Over the fixed orientation-reversing base boundary gluing, it glues the
\emph{entire preimages} of these two circles by the fibre map
\[
 x\longmapsto x a_e,\qquad a_e\in G.
\]
This map is compatible with the two peripheral actions because
$p_{v,j}^{-1}(xa_e)=(p_{u,i}x)a_e$.
It therefore extends uniquely to the required map of boundary covering
spaces.  Every port is used at most once.  All vertex and port labels,
the local monodromy representations, and the underlying base gluing are
fixed when comparing two phase assignments $a=(a_e)$ and $a'=(a'_e)$.

This promise is substantial.  A general permutation of the boundary
lifts, with independent rotations on them, need not be right
multiplication by one group element.  Also, a natural action of a
dihedral group on $m$ points is usually not the regular $2m$-point action
used here.  The existing affine surface-cover classifier contains regular
paired components that can motivate this interface, but no automatic
extraction or conversion is asserted by the new module.

\subsection{From gluing maps to a single root variable}

\begin{lemma}[Local transports]
Every isomorphism of a block cover over the identity of $F_v$ is uniquely
of the form $x\mapsto xh_v$, with $h_v\in G$.
\end{lemma}
\begin{proof}
An equivariant map is determined by $h_v=f(1)$, since
$f(g)=gf(1)=gh_v$.  Conversely, right multiplication commutes with every
left monodromy action and is bijective.  The associated covering-space
isomorphism is unique by path lifting.
\end{proof}

\begin{theorem}[Exact assembly transport equations]
Two phase assignments are isomorphic over the fixed labelled base if and
only if there are $h_v\in G$ satisfying
\begin{equation}\label{eq:regular-glue}
 a_eh_v=h_ua'_e\qquad(e:u\to v).
\end{equation}
For a spanning forest with roots $r$, let $A_v$ and $A'_v$ be the ordered
products of the source and target edge phases along the forest path from
$r$ to $v$, using inverses when an edge is traversed backwards.  Then
every solution is uniquely specified by its root elements $h_r$, with
\begin{equation}\label{eq:regular-root}
 h_v=A_v^{-1}h_rA'_v.
\end{equation}
For each nonforest edge, the remaining condition is
\begin{equation}\label{eq:regular-cycle}
 c_eh_r=h_rc'_e,\qquad
 c_e=A_ua_eA_v^{-1},\quad
 c'_e=A'_ua'_e(A'_v)^{-1}.
\end{equation}
\end{theorem}
\begin{proof}
Starting with $x$ on the source side of an edge, gluing and then applying
the target-block transport gives $xa_eh_v$.  Transporting first and then
using the target gluing gives $xh_ua'_e$.  Equality for all $x$ is exactly
\eqref{eq:regular-glue}.  It is necessary and sufficient for the local
cover isomorphisms to descend to an isomorphism of the glued covers.
The equations on a tree uniquely propagate a chosen root element and
give \eqref{eq:regular-root} by induction.  Substituting that expression
into the equation on a chord and multiplying on the two sides yields
\eqref{eq:regular-cycle}.  The calculation includes loops and reverse
tree traversals.
\end{proof}

For cyclic groups, \eqref{eq:regular-cycle} simply says $c_e=c'_e$:
the root variable disappears.  For $D_m$, write
$h_r=(s,t)$, $c_e=(p,a)$, and $c'_e=(p',a')$.  The exact conditions are
\[
 p=p',\qquad
 \begin{cases}
 a=(-1)^sa'\pmod m,&p=0,\\
 2t=a-(-1)^sa'\pmod m,&p=1.
 \end{cases}
\]
There are only two \emph{global} parity choices per connected graph
component, regardless of the number of blocks.  The second congruence
is impossible when $\gcd(2,m)$ does not divide its right side; otherwise
it is one residue class modulo $m/\gcd(2,m)$.

\subsection{Ordered points and unparameterized boundary lifts}

A boundary lift of a port with peripheral monodromy $p$ is represented
by a sheet $x$.  Its fibre over the boundary basepoint is the left coset
$\langle p\rangle x$.  Different sheets in that coset describe the same
unparameterized boundary circle.  A named-boundary comparison therefore
must preserve a coset, rather than an arbitrarily selected point on it.

\begin{lemma}[Mark reduction]
Suppose that at vertex $v$ a source sheet $x$ is to map to a target sheet
$y$.  In the forest coordinates this is the exact root condition
\[
 h_r=A_vx^{-1}y(A'_v)^{-1}.
\]
If the mark specifies only the boundary lift containing $x$ and the one
containing $y$, then the condition is
\begin{equation}\label{eq:root-coset}
 h_r\in L\langle p\rangle Q,\qquad
 L=A_vx^{-1},\quad Q=y(A'_v)^{-1}.
\end{equation}
For each fixed root parity, this is either impossible or one congruence
$t=b\pmod q$ with $q\mid m$.
\end{lemma}
\begin{proof}
The point statement follows directly from $xh_v=y$.  For boundary lifts,
the required equality is
$\langle p\rangle xh_v=\langle p\rangle y$, which is equivalent to
\eqref{eq:root-coset}.
If $p=T^a$, put $q=\gcd(m,a)$.  The elements of
$L\langle p\rangle Q$ have the parity of $LQ$, and their shifts form the
single residue class of the shift of $LQ$ modulo $q$.  This follows by
multiplying $L(T^{ka})Q$; the possible sign on $ka$ does not change its
generated cyclic subgroup.
If $p=T^aR$, the subgroup has exactly two elements.  The two elements
$LQ$ and $LpQ$ have opposite parities, so each parity permits exactly
one shift modulo $m$.  This also proves the degenerate cases $m=1,2$.
\end{proof}

Repeated marks and an ordered list of boundary names are permitted.
The constraint for each named lift is independent of the representative
sheet used to encode it.  Individual attaching-circle rotations are
not silently forgotten: edge phases remain part of the input, and
parameterized point marks can fix phases that an unparameterized
boundary mark leaves free.

\subsection{The complete binary algorithm and its certificates}

\begin{theorem}[Compressed marked assembly comparison]
For the promised cyclic or dihedral assembly model, with ordered sheet
marks and ordered unparameterized boundary-lift marks, there is a
deterministic algorithm that computes \emph{all} isomorphisms over the
fixed base in polynomial binary time.
For each graph component, its answer is the union of at most two
families
\[
 h_r=T^{b+kq}R^s,\qquad
 0\le k<m/q,\quad q\mid m.
\]
These families, together with \eqref{eq:regular-root}, represent every
compatible map uniquely.  The cyclic case has at most one family.
\end{theorem}
\begin{proof}
Validate the canonical bordered-surface schema, the surjectivity of the
local group representations, distinct use of ports, and inverse peripheral
monodromies.  Construct the spanning forest and the path products.
For each root parity, reduce the cycle and mark conditions as above.
Every nontrivial surviving condition is a single modular congruence
$t=b_i\pmod{q_i}$ with $q_i\mid m$.

Intersect them by generalized Chinese remaindering.  If the current class
is $t=b\pmod q$ and the next is $t=b'\pmod{q'}$, their intersection exists
if and only if
\[
 b=b'\pmod{\gcd(q,q')}.
\]
When it exists, extended Euclid produces the one resulting residue
class modulo $\operatorname{lcm}(q,q')$, still a divisor of $m$.
When it does not, there are two incompatible original congruences:
pairwise consistency of such congruences is equivalent to simultaneous
consistency.  A scan of the already processed constraints locates that
pair; it occurs at most once per parity branch.

The result for a surviving branch is exactly the stated arithmetic
progression.  Root propagation is unique, so no maps are missing or
counted twice.  Different parities are different elements of the regular
dihedral group, including when $m=1,2$.  Different graph components have
independent choices.
\end{proof}

Here is a useful precise complexity accounting.  Let $N$ be the total
number of vertices, edges, free generators in all local schemas, and marks,
and let $B$ include the maximum input-integer bit length and $\log(N+2)$.
For a connected graph, validation, comparison, canonicalization, and
certificate verification have a conservative $O(NB^3)$ elementary
bit-operation bound.  There are $O(N)$ arithmetic operations, with
extended Euclid and modular inversion charged conservatively.  All
residue moduli divide $m$.  The representation contains $O(NB)$ bits.
For a disconnected graph the residue-family computation remains
$O(NB^3)$; explicitly multiplying all component isomorphism counts adds
at most $O(N^2B^2)$ work with schoolbook multiplication.  One may retain
the factored count if that product is unnecessary.

The implementation emits arithmetic certificates.  A positive parity
branch gives a residue satisfying every reduced congruence and the least
common multiple of their moduli.  A negative branch identifies either
one impossible equation or two incompatible congruences.  The verifier
reconstructs the forest and reduced conditions, but does \emph{not} run
CRT.  A second verifier checks an explicit set of block gauges directly
against \eqref{eq:regular-glue} and the original point/coset marks, without
using a forest or a congruence solver.  Thus a successful transport is
checked against its actual gluing semantics.

\subsection{Canonical keys for unmarked assemblies}

\begin{theorem}[Canonical unmarked representative]
For a fixed labelled cyclic or dihedral assembly model, a complete
unmarked isomorphism key can be computed in polynomial binary time.
In the dihedral case at most two candidate root conjugations per graph
component suffice.  The algorithm also produces block gauges to the
canonical edge-phase assignment.
\end{theorem}
\begin{proof}
First gauge the forest phases to the identity.  The remaining phases are
the ordered tuple $(c_e)$ of chord holonomies.  Residual changes are
simultaneous conjugation by one root element.
For a cyclic group there is nothing further to do.

In $D_m$, conjugation by $h=(s,t)$ sends a rotation shift $a$ to
$(-1)^sa$ and a reflection shift $a$ to $(-1)^s(a-2t)$.
If every chord is a rotation, take the lexicographically smaller of
the tuple of shifts and its negation.
Otherwise let $a$ be the shift of the first reflection chord in the fixed
edge order, set $d=\gcd(2,m)$ and $\delta=a\bmod d$, and require that this
reflection have the smallest attainable shift $\delta$.  For each parity
$s$, solve
\[
 2t=a-(-1)^s\delta\pmod m.
\]
The equation always has solutions.  For even $m$ its two solutions differ
by $m/2$, and $T^{m/2}$ is central, so they act identically by conjugation
on every chord.  There is therefore only one relevant tuple per parity.
Choose the lexicographically smaller of these at most two tuples.
The normalized forest and chord phases give a complete canonical phase
assignment.  Include the local schemas, monodromies, and labelled port
graph in the key so it cannot identify different framed input models.
\end{proof}

The cache key is intentionally \emph{unmarked}.  It must not be used to
erase external named lifts, points, base vertex identities, or arbitrary
attachment maps.  The marked comparison algorithm handles named points
and lifts exactly, but a canonical marked-state representation beyond
this scope is a separate task.

\subsection{Exact residual state counts}

\begin{theorem}[Phase orbits and their exact counts]
Let the fixed labelled graph be connected with cycle rank
$\beta=|E|-|V|+1$.  The set of unmarked phase assignments modulo all block
deck gauges is naturally identified with $G^\beta$ modulo simultaneous
conjugacy.  Consequently, for $C_m$ it has exactly $m^\beta$ elements.
For $D_m$ its cardinality is
\[
 \begin{cases}
 \displaystyle
 \frac{(2m)^\beta+(m-1)m^\beta+m\,2^\beta}{2m},
 &m\text{ odd},\\[2ex]
 \displaystyle
 \frac{2(2m)^\beta+(m-2)m^\beta+m\,4^\beta}{2m},
 &m\text{ even}.
 \end{cases}
\]
\end{theorem}
\begin{proof}
Gauge each forest edge to the identity.  Any choice of phases on the
$\beta$ remaining edges is possible, and the only remaining freedom is
the common root gauge acting by simultaneous conjugacy.  This proves the
first statement and the cyclic count.
Burnside's lemma gives the dihedral count in the form
\[
 \frac1{|D_m|}\sum_{g\in D_m}|C_{D_m}(g)|^\beta.
\]
For odd $m$, the identity has centralizer size $2m$, the $m-1$ other
rotations have centralizer size $m$, and each of the $m$ reflections has
centralizer size $2$.  For even $m$, the two central rotations have
centralizer size $2m$, the $m-2$ other rotations have centralizer size $m$,
and the $m$ reflections have centralizer size $4$.  Substitution proves
the displayed formulas.  The same expressions include $m=1,2$ and
$\beta=0$.
\end{proof}

These formulas distinguish cheap equivalence from a small search space.
Even with $\beta=1$, a cyclic modulus $m=2^b$ gives $2^b$ inequivalent
phase assignments encoded with $O(b)$ bits each.  The new algorithm
removes enumeration of the vertex gauges for a \emph{given comparison};
it does not reduce all possible holonomy choices to polynomially many
states.  Equally, this count does not prove that an unknot hierarchy
must encounter all these states.  A quasipolynomial recognizer needs a
bound on the states reachable by its geometric construction, or a
further proved quotient appropriate to the information the construction
needs to retain.

Although every block is regular, a glued cover need not be regular.
Its deck group is the simultaneous centralizer of its chord holonomies
inside $G$ (up to the chosen right-action convention).  In particular,
noncentral dihedral holonomy can reduce the global deck group.  The
algorithm classifies these assemblies without assuming global regularity.

\subsection{Topology and a concrete large example}

Each connected graph component produces a connected assembled covering
surface.  Boundary-circle gluings do not change Euler characteristic.
Thus, writing $d=|G|$,
\[
 \chi(\widetilde F)=d\sum_{v}\chi(F_v),\qquad
 b(\widetilde F)
 =\sum_{\text{unglued ports }(v,j)}
        \frac{d}{\operatorname{ord}(p_{v,j})},
\]
and $g(\widetilde F)=(2-b(\widetilde F)-\chi(\widetilde F))/2$.
For a rotation $T^a$ the order is $m/\gcd(m,a)$; a dihedral reflection
has order two.  These formulas use binary arithmetic and no expanded
boundary-cycle enumeration.

For example, take two pair-of-pants bases with local monodromies $T,R$
in $D_m$.  Their connected regular covers have genus zero and $2m+2$
boundary circles.  Glue the entire preimages of their reflection ports.
There are $m$ glued circle pairs, giving a connected surface of genus
$m-1$ and $2m+4$ remaining boundary circles.  The graph has one edge, so
all gluing phases are equivalent and exactly $2m$ isomorphisms relate
any two assignments.  The construction and comparison still use
polynomial time in $\log m$.  This example shows that exponential
expanded genus is compatible with a small, exact assembly representation.

\subsection{Relevance to an unknot hierarchy and the remaining obligations}

Lackenby's July 2026 hierarchy algorithm explicitly retains boundary
patterns and parallelity information \cite{lackenby2026}.  Its
Section~9 iteration estimate does not by itself give a quasipolynomial
implementation; the treatment of supplied compressed surfaces is a
different issue from controlling which states are generated.
General orbit algorithms for compressed normal surfaces are already
available \cite{aht-orbits}.  The present kernel should be evaluated as
an additional exact attachment-comparison tool, not as a replacement
for those algorithms.

The next geometric work is concrete:
\begin{enumerate}
\item Extract regular cyclic or dihedral block actions from a supplied
compressed surface or parallelity description, with certificates for
the common fibre coordinates and the based peripheral system.
\item Certify that each proposed full-boundary-preimage attachment is
one coherent right-group phase.  If the actual map is an arbitrary
permutation of lifts with independent rotations, this module does not
apply.
\item Determine whether the coherent-phase promise is preserved by the
specific cutting, compression, and hierarchy-repair operations used in
the recognizer.
\item Extend the interface to controlled exceptions, such as a bounded
number of independently attached lifts or nonregular blocks, while
retaining exact external marking semantics.
\item Bound the cycle rank and the reachable holonomy tuples created
by the hierarchy, rather than only the time needed to compare two
tuples.
\item Investigate whether a weaker equivalence than cover isomorphism
preserves all later pattern-compression decisions and permits additional
state merging; prove the preservation theorem before using that quotient.
\end{enumerate}

Until those obligations are met, the module returns no knot verdict and
is not dispatched by the main recognition pipeline.  Its proved
performance improvement is the removal of sheet and local-gauge
enumeration from this explicit geometric assembly problem.
